
We asked whether reproducibility badges improve reproducibility outcomes in machine learning. Our answer: badges successfully increase artifact presence but do not ensure artifact quality, producing no detectable variance reduction across independent reproduction attempts.

This study provides the first continuous quality-outcome measurement linking documentation artifact quality to reproducibility in ML benchmarks. Analyzing 108 classification benchmarks (2019-2024) with 5+ independent reproduction attempts each, we found mean artifact quality of 2.43/10, far below the threshold for actionable implementation guidance (7.0/10). Critical dimensions---evaluation protocols (1.19/10) and hyperparameters (1.16/10)---were almost never documented. Automated scoring consistency $\kappa$=1.0 confirms this reflects genuine quality deficits, not measurement noise.

Despite artifact deposition at scale, performance variance (CV) showed no statistically significant reduction: high-artifact benchmarks (mean CV=0.035) versus low-artifact benchmarks (mean CV=0.069), Mann-Whitney p=0.418, Cohen's d=0.464. While the directional trend is positive, the effect size falls below the medium threshold (0.5) and the sample (n=22 benchmarks in variance analysis) was underpowered relative to our target (n=100, 80\% power). Our findings suggest reproducibility badge programs face a pattern consistent with checkbox compliance: authors create artifacts to satisfy venue requirements but do not invest in documentation quality.

The main conclusion is clear: artifact presence is insufficient---quality enforcement is needed. Current badge programs incentivize deposition (GitHub repositories, dataset cards) but lack post-publication quality audits. Venues should complement presence incentives with quality mechanisms: rubric-based scoring as a submission requirement, post-publication community review, or automated completeness checks for evaluation protocols and hyperparameters.

Future work should expand the sample to n=100 benchmarks to achieve adequate statistical power, enabling definitive conclusions about the weak effect observed (d=0.464). Quality-weighted scoring---using our rubric to predict variance reduction as a continuous function of artifact quality rather than binary presence---would test whether high-quality artifacts succeed where low-quality ones fail. Stratification by venue prestige (top-tier versus mid-tier conferences) could reveal whether enforcement norms differ across communities. Longitudinal analysis of artifact decay (link rot, repository archival) would measure whether quality deteriorates post-publication. Finally, automated artifact quality assessment tools could enable real-time feedback during submission, shifting the burden from post-hoc auditing to proactive guidance.
